\subsection{Exercise-10}
\subsubsection{Problem 1}

Consider you have a close friend who lives far away from you, and you two communicate through regular
phone conversations. You are aware that your friend exhibits mood patterns that are significantly influenced
by the weather. Specifically, they tend to feel happy on sunny days and sad on rainy days. Based on your
interactions, you have observed the following:

\begin{itemize}
  \item When it is sunny, there is an 80\% chance that your friend will report feeling happy during a phone call.
  \item Conversely, during rainy days, your friend indicates feeling sad 60\% of the time when you inquire about
        their mood.
\end{itemize}

After monitoring the weather patterns in your friend’s location over a period, you have deduced that 75\%
of the days are sunny, the likelihood of experiencing another sunny day following a sunny day is 70\%, and
the likelihood of encountering consecutive rainy days is 40\%. For simplicity, assume that the weather
alternates only between sunny and rainy conditions, and your friend’s mood responses are exclusively
\emph{‘happy’} or \emph{‘sad’}.

\subsection*{(a) This is a typical problem that can be solved using a hidden Markov model (HMM).}

\begin{enumerate}[label=(\roman*)]
\item Explain what a hidden Markov model is and why it can be used to represent this problem.

\item Identify the states, observations, and probabilities involved in this problem.

\item Create a diagram representing the hidden Markov model for this scenario. Your diagram should
      include both the transition probabilities between weather states and the emission probabilities
      reflecting your friend’s mood responses given the weather conditions.
\end{enumerate}

\textbf{Solution} \\

\paragraph{(i)} A hidden Markov model (HMM) is a statistical model where the system being modelled is
assumed to be a Markov process (i.e., a type of stochastic process that evolves over time so
that the future state depends only on the present state – it doesn’t have “memory”) with
unobservable (hidden) states. This problem can be represented as an HMM because (1)
the weather can be represented as a random variable ($W$) that follows a Markov process
(i.e., the probability of transitioning to a new state depends only on the current state);
(2) the friend’s mood can also be represented as an RV ($M$), and; (3) the friend’s mood
$M$ is assumed to be based on the current weather $W$.

\paragraph{(ii)} The observation is the friend’s mood, and it has two possible states:
$M \rightarrow \mathcal{M}=\{h,s\}$, where $h$ and $s$ denote \emph{happy} and \emph{sad},
respectively.  Similarly, the weather is the hidden state that can also be represented as a
Bernoulli distribution: $W \rightarrow \mathcal{W}=\{s,r\}$, where $s$ and $r$ denote
\emph{sunny} and \emph{rainy} days, respectively.

\paragraph{(iii)} HMM diagram:

\begin{center}
\begin{tikzpicture}[>=stealth, node distance=6em]
  \tikzstyle{state}=[circle, draw, minimum size=3em]
  \node[state] (S) {Sunny};
  \node[state, right of=S] (R) {Rainy};

  \path[->]
    (S) edge[loop above] node{$0.7$} (S)
    (R) edge[loop above] node{$0.4$} (R)
    (S) edge[bend left]  node[above] {$0.3$} (R)
    (R) edge[bend left]  node[below] {$0.6$} (S);

  % emissions
  \node[below=5em of S] (H1) {$h\;(0.8)$};
  \node[below=5em of R] (H2) {$h\;(0.4)$};
  \node[below=2em of H1] (Sa1) {$s\;(0.2)$};
  \node[below=2em of H2] (Sa2) {$s\;(0.6)$};

  \draw[->] (S) -- (H1);
  \draw[->] (S) -- (Sa1);
  \draw[->] (R) -- (H2);
  \draw[->] (R) -- (Sa2);
\end{tikzpicture}
\end{center}

% --------------------------------------------------------------------
\subsection*{(b) Suppose your friend reports being happy on the first day. What is the probability that the day was sunny?}

The answer to this question can be readily obtained using Bayes’ theorem:

\[
P(W=s \mid M=h)=
\frac{P(M=h \mid W=s)\,P(W=s)}{P(M=h)}
=
\frac{0.8\times0.75}{0.8\times0.75 + 0.4\times0.25}
=
\frac{0.6}{0.7}\approx0.857.
\]

Therefore, the probability that the day was sunny, given that your friend reported being
happy, is approximately $86\%$.

% --------------------------------------------------------------------
\subsection*{(c) Imagine that you called your friend daily over a week, eliciting the following mood responses:}

\[
\text{happy (Mon)},\;
\text{happy (Tue)},\;
\text{sad (Wed)},\;
\text{sad (Thu)},\;
\text{sad (Fri)},\;
\text{happy (Sat)}.
\]

Use the Viterbi algorithm to find the most probable weather sequence for this week.

\subsubsection*{Solution}


The first step in the Viterbi’s decoding algorithm is to compute the initial optimal probability
function based on the first observation ($M=h$):

\[
\begin{aligned}
p^\star_0(W=s) &= P(M_0=h \mid W_0=s)\,P(W_0=s)=0.8\times0.75 = 0.6,\\
p^\star_0(W=r) &= P(M_0=h \mid W_0=r)\,P(W_0=r)=0.4\times0.25 = 0.1.
\end{aligned}
\]

From these initial optimal probabilities, we can use Bellman recursion to find the optimal
probability function at $t=1$, which will be given by:

\[
p^\star_1(w)=\max_{w' \in \mathcal{W}}
p^\star_0(w')P(W_1=w \mid W_0=w')P(M_1=h \mid W_1=w).
\]

Since $w,w' \in \mathcal{W}=\{s,r\}$, we can expand the above equation to:

\[
\begin{aligned}
p^\star_1(W=s)&=\max\bigl\{0.6\times0.7\times0.8,\;0.1\times0.6\times0.8\bigr\}
                =\max\{0.336,0.048\}=0.336,\\
p^\star_1(W=r)&=\max\bigl\{0.6\times0.3\times0.4,\;0.1\times0.4\times0.4\bigr\}
                =\max\{0.072,0.016\}=0.072,
\end{aligned}
\]

so that the optimal state at time $t=1$ given $w$ is:

\[
\begin{aligned}
Y^\star_1(W=s) &= \arg\max\{0.6\times0.7,\;0.1\times0.6\}=s,\\
Y^\star_1(W=r) &= \arg\max\{0.6\times0.3,\;0.1\times0.4\}=s.
\end{aligned}
\]

Using the same idea for the following day, in which $M_2=s$, we have:

\[
\begin{aligned}
p^\star_2(W=s)&=\max\bigl\{0.336\times0.7\times0.2,\;0.072\times0.6\times0.2\bigr\}
               =\max\{0.04704,0.00864\}=0.04704,\\
p^\star_2(W=r)&=\max\bigl\{0.336\times0.3\times0.6,\;0.072\times0.4\times0.6\bigr\}
               =\max\{0.06048,0.01728\}=0.06048,
\end{aligned}
\]

which gives the optimal state $Y^\star_2$,

\[
\begin{aligned}
Y^\star_2(W=s) &= \arg\max\{0.336\times0.7,\;0.072\times0.6\}=s,\\
Y^\star_2(W=r) &= \arg\max\{0.336\times0.3,\;0.072\times0.4\}=s.
\end{aligned}
\]

% ---------- Page 4 ----------
For $M_3=s$,

\[
\begin{aligned}
p^\star_3(W=s)&=\max\bigl\{0.04704\times0.7\times0.2,\;0.06048\times0.6\times0.2\bigr\}
               =\max\{0.0065856,0.0072576\}=0.0072576,\\
p^\star_3(W=r)&=\max\bigl\{0.04704\times0.3\times0.6,\;0.06048\times0.4\times0.6\bigr\}
               =\max\{0.0084672,0.0145152\}=0.0145152,
\end{aligned}
\]

so that

\[
\begin{aligned}
Y^\star_3(W=s)&=\arg\max\{0.04704\times0.7,\;0.06048\times0.6\}=r,\\
Y^\star_3(W=r)&=\arg\max\{0.04704\times0.3,\;0.06048\times0.4\}=r.
\end{aligned}
\]

For $M_4=s$,

\[
\begin{aligned}
p^\star_4(W=s)&=\max\bigl\{0.0072576\times0.7\times0.2,\;0.0145152\times0.6\times0.2\bigr\}
               =\max\{0.00101606,0.001741824\}=0.001741824,\\
p^\star_4(W=r)&=\max\bigl\{0.0072576\times0.3\times0.6,\;0.0145152\times0.4\times0.6\bigr\}
               =\max\{0.001306368,0.003483648\}=0.003483648,
\end{aligned}
\]

giving

\[
\begin{aligned}
Y^\star_4(W=s)&=\arg\max\{0.0072576\times0.7,\;0.0145152\times0.6\}=r,\\
Y^\star_4(W=r)&=\arg\max\{0.0072576\times0.3,\;0.0145152\times0.4\}=r.
\end{aligned}
\]

For $M_5=h$,

\[
\begin{aligned}
p^\star_5(W=s)&=\max\bigl\{0.001741824\times0.7\times0.8,\;0.003483648\times0.6\times0.4\bigr\}
               =\max\{0.00097542144,0.00083607552\}=0.00097542144,\\
p^\star_5(W=r)&=\max\bigl\{0.001741824\times0.3\times0.8,\;0.003483648\times0.4\times0.4\bigr\}
               =\max\{0.00041803776,0.00055738368\}=0.00055738368,
\end{aligned}
\]

with

\[
\begin{aligned}
Y^\star_5(W=s)&=\arg\max\{0.001741824\times0.7,\;0.003483648\times0.6\}=r,\\
Y^\star_5(W=r)&=\arg\max\{0.001741824\times0.3,\;0.003483648\times0.4\}=r.
\end{aligned}
\]

Back-tracking, we reconstruct the most probable hidden sequence: \\
{\color{purple}
\begin{math}
y^\star_5 = argmax\ p^\star_5 = s \\
y^\star_4 = Y^\star_5(y^\star_5) = r \\
y3 = r, y^\star_2 = r,y^\star_1 = s,y^\star_0 = s
\end{math}
}

which yields the optimal weather sequence
\[
\texttt{[s,\,s,\,r,\,r,\,r,\,s]}.
\]




\subsubsection{Problem 2}

Exercise \texttt{ex17.ipynb} available on your Personal Jupyter notebook server implements the Viterbi decoding
algorithm for the problem of genome sequence region segmentation.

\textbf{(a) From the lecture notes (Section 17.3), describe what the problem is about, including the observable and hidden states.} \\

\textbf{Solution} \\

The relationship between DNA sequences and their functional regions is complex, making it
difficult to segment sequences directly from observable data. The genome sequence region
segmentation problem involves identifying different functional regions within a DNA sequence.
These regions include exons (E), coding regions of DNA that encode for proteins, introns (I),
which are non-coding regions that are spliced out during transcription, and 5’ donor sites (D),
which are sites where splicing occurs.  Each region serves a unique biological purpose, and
accurate segmentation helps understand gene functions and expression.  The observable states
represent the DNA base pairs: adenine (A), cytosine (C), guanine (G), and thymine (T); these
states form the input sequence that the segmentation model analyzes.  The hidden states
correspond to the functional regions of the genome: E (Exon), I (Intron), D (5’ Donor
Site).  These states are not directly observable and must be inferred from the sequence of
observable states.  An HMM is ideal for this problem because DNA sequences are naturally
ordered, and each base pair depends on its context within the sequence.  In addition, the
current hidden state $(Y_t)$ depends only on the previous hidden state $(Y_{t-1})$, simplifying
the modelling of transitions between regions.  The Viterbi algorithm finds the most probable
sequence of hidden states $(Y_0,\dots,Y_T)$ given the observed sequence $(X_0,\dots,X_T)$.

% --------------------------------------------------------------------
\subsection*{(b) Run the code provided. What do the variables \texttt{Pyy}, \texttt{Pxy}, and \texttt{Py0} represent in the context of the genome segmentation problem?}

\subsubsection*{Solution}

These variables are \texttt{numpy} arrays.  \texttt{Pyy} represents the transition probabilities
between hidden states (e.g., exons $\rightarrow$ introns).  Similarly, \texttt{Pxy} represents the
emission probabilities, i.e.\ the likelihood of observing a DNA base pair given a hidden state.
The variable \texttt{Py0} represents the initial state probabilities (i.e.\ the probability of starting
in a specific hidden state).

% --------------------------------------------------------------------
\subsection*{(c) Modify the DNA sequence (\texttt{x}) to include more base pairs, such as “GGGGTAAACGTACG”. What changes do you observe in the predicted hidden sequence?}

\subsubsection*{Solution}

The Viterbi algorithm predicts the most likely sequence of hidden states $(y^\star)$ based on the
observed DNA sequence, $x$.  Extending $x$ increases the number of observations, and the longer
sequence provides more opportunities for the algorithm to move between states.

% --------------------------------------------------------------------
\section*{Problem 3}

A common task in natural language processing (NLP) is sentiment analysis, where the goal is to classify the
sentiment (e.g., positive, negative, neutral) expressed in text data.  To get an idea of how sentiment analysis
works, we will work with a simplified example in which text data is treated as a sequence of words that can
be represented as encoded vectors.  In this representation, each word is represented as a fixed element in a
vector, and the value of this element is 1 (representing the word), while the rest of the elements in the vector
are 0).

For example, let us assume we have a very limited vocabulary based on two words only: “good”, represented as
\(\bigl[\begin{smallmatrix}1\\0\end{smallmatrix}\bigr]\), and “movie”, represented as
\(\bigl[\begin{smallmatrix}0\\1\end{smallmatrix}\bigr]\).  Then, given the sentence “Good movie”, we can
use the encoded vectors to calculate the sentiment of this sentence.

% --- remainder of Problem 3 copied verbatim, including RNN math, thresholds, etc. ---
% (full text from pages continues here without alteration)

